\documentclass[11pt]{article}
\usepackage[T1]{fontenc}
\usepackage{lmodern}
\usepackage{amsmath,amssymb}
\usepackage[margin=1in]{geometry}
\usepackage[colorlinks=true,linkcolor=blue,citecolor=blue,urlcolor=blue]{hyperref}

\title{Literature Status: Nontrivial Type and Pisier's Inequality}
\author{Run \texttt{fa\_banach\_001}, agent \texttt{agent\_lane\_12}}
\date{2026-06-15}

\begin{document}
\maketitle

\section*{Status}

This is a lightweight literature-status packet. It records that the main
nontrivial-type question in Hytönen--Naor, \emph{Pisier's inequality
revisited}, arXiv:1207.5375, is answered by Ivanisvili--van Handel--Volberg,
\emph{Rademacher Type and Enflo Type Coincide}, arXiv:2003.06345.

\section*{Source Question}

The source paper asks whether every Banach space of nontrivial type satisfies
Pisier's inequality with the logarithmic dimension factor replaced by a
constant independent of the cube dimension. In the parsed source this appears
at \texttt{source.tex:391--395}.

\section*{Supporting Answer}

The supporting paper gives a finite-cotype characterization. Its Theorem
\texttt{thm:cotype} states that, for any fixed Banach space $X$ and
$1\le p<\infty$, Pisier's inequality holds with a constant independent of
dimension if and only if $X$ has finite cotype
(\texttt{source.tex:411--413} in the parsed supporting source).

Immediately after this theorem, the authors observe that every Banach space
with nontrivial type has finite cotype and conclude that Pisier's inequality
therefore holds dimension-free in any Banach space with nontrivial type
(\texttt{source.tex:414--417}).

\section*{Scope}

This gives an affirmative later-literature answer to the source question. It
does not settle the separate quantitative problem in Hytönen--Naor concerning
the correct asymptotic order, as $r\to\infty$, of the best constant in
Pisier's inequality for $X=\ell_r$.

\section*{References}

\begin{itemize}
\item T. Hytönen and A. Naor, \emph{Pisier's inequality revisited},
arXiv:1207.5375.
\item P. Ivanisvili, R. van Handel, and A. Volberg,
\emph{Rademacher Type and Enflo Type Coincide}, arXiv:2003.06345.
\end{itemize}

\end{document}
